% Preamble
\documentclass[11pt, aspectratio=169]{beamer}
	% Packages
	\usepackage[english]{babel}
	\usepackage{lipsum}
	\usepackage{mathptmx}
	\usepackage{xcolor}
	% Themes
	\usefonttheme[onlymath]{serif}
	\usetheme[nofirafonts]{focus}
	\renewcommand{\thempfootnote}{\arabic{mpfootnote}}
	% Cover
	\title{Microeconometrics Using Stata}
	\subtitle{Linear regression basics: Exercises}
	\author{
		Jhon R. Ordoñez
		\footnote{\label{1} National University of San Cristóbal de Huamanga}
		\footnote{\label{2} Faculty of Economic, Administrative and Accounting Sciences}
		\footnote{\label{3} Professional School of Economics}
	}
	\date{\today}
% Body
\begin{document}
	% Cover
	\begin{frame}
		\maketitle
	\end{frame}
	% Outline
	\begin{frame}[t]
		\frametitle{Outline}
		\tableofcontents
	\end{frame}
	% Section 1
	\section{Exercise 1}
		% Slide 1.1
		\begin{frame}
			\frametitle{Exercise 1}
				\begin{block}{Exercise 1}
					Fit the model in \textcolor{red!50}{\textbf{section 3.5}} using only the first $100$ observations.
					Compute standard errors in three ways: default, heteroskedastic, and
					cluster-robust where clustering is on the number of chronic problems.
					Use \textcolor{blue}{\textbf{estimates}} to produce a table with three sets of coefficients and
					standard errors, and comment on any appreciable differences in the
					standard errors. Construct a similar table for three alternative sets of
					heteroskedasticity-robust standard errors, obtained by using the
					\textcolor{blue}{\textbf{vce(robust)}}, \textcolor{blue}{\textbf{vce(hc2)}}, and \textcolor{blue}{\textbf{vce(hc3)}} options, and comment on any
					differences between the different estimates of the standard errors.
				\end{block}
		\end{frame}
	% Section 2
	\section{Exercise 2}
		% Slide 2.1
		\begin{frame}
			\frametitle{Exercise 2}
			\begin{block}{Exercise 2}
				Fit the model in \textcolor{red!50}{\textbf{section 3.5}} with robust standard errors reported. Test
				at $5\%$ the joint significance of the demographic variables \textcolor{blue}{\textbf{age}}, \textcolor{blue}{\textbf{female}},
				and \textcolor{blue}{\textbf{income}}. Test the hypothesis that being male (rather than female)
				has the same impact on medical expenditures as aging 10 years. Fit the
				model under the constraint that $\beta_{phylim} = \beta_{actlim}$ by first typing
				\textcolor{blue}{\textbf{constraint 1 phylim = actlim}} and then by using \textcolor{blue}{\textbf{cnsreg}} with the
				\textcolor{blue}{\textbf{constraints(1)}} option.
			\end{block}
		\end{frame}
	% Section 3
	\section{Exercise 3}
		% Slide 3.1
		\begin{frame}
			\frametitle{Exercise 3}
			\begin{block}{Exercise 3}
				Fit the model in \textcolor{red!50}{\textbf{section 3.6}}, and implement the RESET test manually by
				regressing $y$ on $x$ and $\hat{y}^{2}$, $\hat{y}^{3}$, and $\hat{y}^{4}$ and jointly testing that the
				coefficients of $\hat{y}^{2}$, $\hat{y}^{3}$, and $\hat{y}^{4}$ are $0$. 
				To get the same results as \textcolor{blue}{\textbf{estat ovtest}}, 
				do you need to use default or robust estimates of the \textbf{VCE} in
				this regression? Comment. Similarly, implement \textcolor{blue}{\textbf{linktest}} by
				regressing $y$ on $\hat{y}$ and $\hat{y}^{2}$ and testing that the coefficient of $\hat{y}^{2}$ is $0$. To
				get the same results as \textcolor{blue}{\textbf{linktest}}, do you need to use default or robust
				estimates of the \textbf{VCE} in this regression? Comment.
			\end{block}
		\end{frame}
	% Section 4
	\section{Exercise 4}
		% Slide 4.1
		\begin{frame}
			\frametitle{Exercise 4}
			\begin{block}{Exercise 4}
				Fit the model in \textcolor{red!50}{section 3.6}, and perform the standard Lagrange
				multiplier test for heteroskedasticity by using \textcolor{blue}{\textbf{estat hettest}} with
				$\mathbf{z} = \mathbf{x}$. Then, implement the test manually as $0.5$ times the explained
				sum of squares from the regression of $y_{i}^{*}$ on an intercept and $z_{i}$, where
					\begin{equation}
						y_{i}^{*} = \left\{\dfrac{\hat{u}_{i}^{2}}{\dfrac{1}{N}\sum_{j} \hat{u}_{j}^{2}  } \right\} - 1
					\end{equation}
				and $\hat{u}_{i}$ is the residual from the original \textbf{OLS regression}. Next, use \textcolor{blue}{\textbf{estat hettest}} with the iid option, and
				show that this test is obtained as  $N\times R^{2}$, where $R^{2}$ is obtained from
				the regression of $\hat{u}_{i}^{2}$ on an intercept and $\mathbf{z}_{i}$.
			\end{block}
		\end{frame}
	% Section 5
	\section{Exercise 5}
		% Slide 5.1
		\begin{frame}
			\frametitle{Exercise 5}
			\begin{block}{Exercise 5}
				Using the \textbf{DGP} of \textcolor{red!50}{\textbf{section 3.7.6}}, generate a sample of size $100$. Regress 
				$y$ on $x$ and an intercept. Apply the \textbf{RESET} omitted variable test. Does
				the test indicate functional form misspecification? Next, use the \textcolor{blue}{\textbf{estat}}
				\textcolor{blue}{\textbf{imtest}} postestimation command to generate \textbf{IM} omnibus diagnostic
				test statistics. If you had applied only this test, what changes to your
				model specification would you make? Would you make additional
				tests before changing the model specification?
			\end{block}
		\end{frame}
	% Section 6
	\section{Exercise 6}
		% Slide 6.1
		\begin{frame}
			\frametitle{Exercise 6}
			\begin{block}{Exercise 6}
				One type of data transformation that is sometimes used in linear
				regression, though not much in econometrics, is variable
				standardization. To standardize a variable you subtract the sample
				mean of that variable from each observation and divide the result by
				the sample standard deviation of that variable. The resulting variable
				with mean zero and standard deviation one is independent of the
				measurement scale of original data. A standardized regression is a
				regression in which the dependent and regressor variables are all
				standardized before running the regression. The standardized
				regression coefficients are interpreted as measuring the effect of one
				standard deviation change in the regressor and hence are in a sense
				comparable across variables. Note that when a regressor is discrete,
				interpreting its impact in units of standard deviation is not natural.
				Reestimate the regression equation of \textcolor{red!50}{\textbf{section 3.5.2}} after standardizing
				all variables. Are the insignificant coefficients the same as those in the
				previous regression? Which variable has the largest coefficient? Does
				that mean it is the most important variable in the regression?
			\end{block}
		\end{frame}
	% References
	\begin{frame}[t,allowframebreaks]
		\frametitle{References}
		\bibliography{biblio.bib}
		% Style
		\bibliographystyle{apalike}
		% Authors
		\nocite{cameron2022-I}
		\nocite{cameron2022-II}
	\end{frame}
	% Cover
	\begin{frame}
		\maketitle
	\end{frame}
\end{document}